\documentclass{techbrief}

\title{Superpositions correspond to multiple decompositions}
\author{Gabriele Carcassi for the AoP collaboration}
\date{}

\begin{document}

\maketitle
	
\begin{abstract}
	We show that the possibility of superpositions in quantum mechanics coincides with the possibility of multiple decompositions of statistical ensembles into pure states.
\end{abstract}
	
\section{Introduction}

Given that the addition of state vectors in quantum mechanics does not represent a physically realizable operation, it is difficult to give quantum superpositions a direct, physically meaningful interpretation. Here we show that the existence of superpositions is equivalent to the existence of a statistical ensemble that can be prepared as a statistical mixture of different pure states, which is a physically meaningful operation. Given that mixtures of different classical states always yields a different ensemble, this insight also tells us why classical superpositions are not possible.

\section{Premises and conditions}

Let us first describe the relevant premises and conditions. The base premise is that we are working with quantum states with their standard representation in terms of complex vector spaces.

\begin{condition}[QST]\label{QST}
	The state of a quantum system is represented by a ray of a complex inner product space $\mathcal{H}$ associated to that system. An ensemble is represented by a density operator $\rho : \mathcal{H} \to \mathcal{H}$, that is a positive semi-definite trace one self-adjoint operator.
\end{condition}

The first condition states that one state can be expressed as a superposition of other two.

\begin{condition}[QSUP-PUR]\label{QSUP-PUR}
	A pure state $\phi$ can be expressed as a superposition of $\psi_1$ and $\psi_2$. That is, $|\phi\> = c_1 |\psi_1\> + c_2 | \psi_2\>$.
\end{condition}

The second condition states that we can find a mixed state $\rho$ that can be expressed as a mixture of $\psi_1$ and $\psi_2$ as well as a mixture of $\phi$ and something else.

\begin{condition}[QSUP-MIX]\label{QSUP-MIX}
	There exists a mixed state $\rho$ that is mixture of $\psi_1$ and $\psi_2$ and also a mixture of $\phi$ and some other states. That is, $\rho = p_{1} \rho_{\psi_1} + p_{2} \rho_{\psi_2} = p_{\phi} \rho_{\phi} + p_{\hat{\phi}} \rho_{\hat{\phi}}$.
\end{condition}

\section{Claim}

The claim is that the two conditions are equivalent over the base premise.

\begin{claim}
	Conditions \ref{QSUP-PUR} and \ref{QSUP-MIX} are equivalent over \ref{QST}.
\end{claim}

\begin{demonstration}
	 Let us show that \ref{QSUP-PUR} implies \ref{QSUP-MIX}. Let $|\phi\> = c_1 |\psi_1\> + c_2 |\psi_2\>$. We define $|\hat{\phi}\> = c_1 |\psi_1\> - c_2 |\psi_2\>$. We have:
	\begin{equation}
		\begin{aligned}
			| \phi \> \< \phi | &= c_1^* c_1 | \psi_1 \> \< \psi_1 | + c_1^* c_2 | \psi_1 \> \< \psi_2 | + c_2^* c_1 | \psi_2 \> \< \psi_1 | + c_2^* c_2 | \psi_2 \> \< \psi_2 | \\
			| \hat{\phi} \> \< \hat{\phi} | &= c_1^* c_1 | \psi_1 \> \< \psi_1 | - c_1^* c_2 | \psi_1 \> \< \psi_2 | - c_2^* c_1 | \psi_2 \> \< \psi_1 | + c_2^* c_2 | \psi_2 \> \< \psi_2 | \\
			| \phi \> \< \phi | &+ | \hat{\phi} \> \< \hat{\phi} | = 2 |c_1|^2 | \psi_1 \> \< \psi_1 | + 2 |c_2|^2 | \psi_2 \> \< \psi_2 |
		\end{aligned}
	\end{equation}\begin{equation}
		\begin{aligned}
			p_{\phi} &= \frac{\<\phi|\phi\>}{\<\phi|\phi\> + \<\hat{\phi}|\hat{\phi}\>} 
			&p_{\hat{\phi}} = \frac{\<\hat{\phi}|\hat{\phi}\>}{\<\phi|\phi\> + \<\hat{\phi}|\hat{\phi}\>} \\
			p_{1} &= \frac{2 |c_1|^2 \< \psi_1 | \psi_1 \> |}{\<\phi|\phi\> + \<\hat{\phi}|\hat{\phi}\>} 
			&p_{2} = \frac{2 |c_2|^2 \< \psi_2 | \psi_2 \> |}{\<\phi|\phi\> + \<\hat{\phi}|\hat{\phi}\>}
		\end{aligned}
	\end{equation}
	\begin{equation}
		\begin{aligned}
			\rho = p_{\phi} \rho_{\phi} + p_{\hat{\phi}} \rho_{\hat{\phi}} &= \frac{\<\phi|\phi\>}{\<\phi|\phi\> + \<\hat{\phi}|\hat{\phi}\>} \frac{| \phi \> \< \phi |}{\< \phi | \phi \>} + \frac{\<\hat{\phi}|\hat{\phi}\>}{\<\phi|\phi\> + \<\hat{\phi}|\hat{\phi}\>} \frac{| \hat{\phi} \> \< \hat{\phi} |}{\< \hat{\phi} | \hat{\phi} \>} \\
			&=\frac{1}{\<\phi|\phi\> + \<\hat{\phi}|\hat{\phi}\>} \left(| \phi \> \< \phi | + | \hat{\phi} \> \< \hat{\phi} |\right) \\
			&=\frac{1}{\<\phi|\phi\> + \<\hat{\phi}|\hat{\phi}\>} \left(2 |c_1|^2 | \psi_1 \> \< \psi_1 | + 2 |c_2|^2 | \psi_2 \> \< \psi_2 |\right) \\
			&=\frac{2 |c_1|^2 \< \psi_1 | \psi_1 \> |}{\<\phi|\phi\> + \<\hat{\phi}|\hat{\phi}\>} \frac{| \psi_1 \> \< \psi_1 |}{\< \psi_1 | \psi_1 \>} + \frac{2 |c_2|^2 \< \psi_2 | \psi_2 \> |}{\<\phi|\phi\> + \<\hat{\phi}|\hat{\phi}\>} \frac{| \psi_2 \> \< \psi_2 |}{\< \psi_2 | \psi_2 \>} \\
			&=p_{1} \rho_{\psi_1} + p_{2} \rho_{\psi_2} \\
		\end{aligned}
	\end{equation}
	Note that $p_{\phi} + p_{\hat{\phi}}=1$, which means $\rho$ is a statistical mixture with the given probability. This also means that $p_1 + p_2 = 1$ as well. We therefore have a mixed state that can be expressed either as a mixture of $\psi_1$ and $\psi_2$ or of $\phi$ and $\hat{\phi}$.
	
	Let us now show that \ref{QSUP-MIX} implies \ref{QSUP-PUR}. Suppose that $\rho = p_{\phi} \rho_{\phi} + p_{\hat{\phi}} \rho_{\hat{\phi}} = p_{1} \rho_{\psi_1} + p_{2} \rho_{\psi_2}$. Since $\rho$ is the mixture of two pure states, it is supported by a two dimensional subspace. Since $\rho$ is a mixture of $\rho_{\psi_1}$ and $\rho_{\psi_2}$, the subspace is the span of $\psi_1$ and $\psi_2$. Since $\phi$ must be in that subspace as well, it must be a superposition of $\psi_1$ and $\psi_2$.
	
	Combining the two results, $\phi$ is a superposition of $\psi_1$ and $\psi_2$ if and only if there is mixed state $\rho$ that can be expressed as a mixture of both $\psi_1$ and $\psi_2$ as well as a mixture of $\phi$ and some other state.
\end{demonstration}

\section{Conclusion}

Reducing pure state superpositions to properties of statistical mixtures means reducing an abstract mathematical property to a physically, operationally, well-defined one. Therefore, it can provide more insights into the physics described by the math.
	
\end{document}